\chapter{Đổi Giờ, Không Đổi Bài}
\chapterintro{Đổi hình thức để giữ nội dung}{Lớp chán không phải lúc nào cũng vì bài dở; nhiều khi chỉ vì cách vào bài quá một màu.}

\section{Đổi người nói trong một phút}
\begin{tinhhuong}{Tình huống}
Thầy cô nói liên tục đã khá lâu. Lớp không hỏng nhưng bắt đầu trôi vì vai trò người học quá thụ động.
\end{tinhhuong}
\begin{goiydungngay}
Nói: \textit{``Bây giờ đổi vai 1 phút. Em nào tóm lại giúp thầy cô bằng ngôn ngữ của lớp mình?''}
\begin{itemize}
\item Chọn 1 đến 2 em là đủ.
\item Không yêu cầu nói đúng sách tuyệt đối.
\item Sau đó thầy cô chốt lại thật gọn.
\end{itemize}
\end{goiydungngay}
\begin{cauchot}
Đổi người nói một chút có thể giữ bài đi tiếp rất lâu.
\end{cauchot}
\ngancach

\section{Đổi tốc độ, không đổi mục tiêu}
\begin{tinhhuong}{Tình huống}
Lớp bắt đầu hụt hơi vì tiết học đang đi đều đều quá lâu.
\end{tinhhuong}
\begin{goiydungngay}
Thông báo: \textit{``30 giây tốc độ cao: mỗi bàn nói ra đúng một ý, không dài dòng.''}
\begin{itemize}
\item Đổi tiết tấu để lớp tỉnh lại.
\item Mỗi bàn chỉ nói một câu.
\item Sau đó quay lại nhịp bình thường.
\end{itemize}
\end{goiydungngay}
\begin{cauchot}
Muốn cứu nhịp, đôi khi chỉ cần đổi tốc độ chứ chưa cần đổi bài.
\end{cauchot}
\ngancach

\section{Đổi cách ghi bài}
\begin{tinhhuong}{Tình huống}
Lớp mệt vì nghe và chép cùng một kiểu quá lâu.
\end{tinhhuong}
\begin{goiydungngay}
Cho học sinh ghi bằng một trong ba cách: 3 từ khóa, 1 câu chốt, hoặc 1 sơ đồ mini.
\begin{itemize}
\item Không bắt tất cả ghi cùng kiểu.
\item Cho 2 em khoe cách ghi khác nhau.
\item Từ đó cả lớp tỉnh lại khá nhanh.
\end{itemize}
\end{goiydungngay}
\begin{cauchot}
Đổi cách ghi cũng là một cách cứu não đang mệt.
\end{cauchot}
\ngancach

\section{Đổi người viết bảng}
\begin{tinhhuong}{Tình huống}
Lớp chậm lại vì mọi ánh nhìn đang dồn về một hướng quen thuộc quá lâu.
\end{tinhhuong}
\begin{goiydungngay}
Mời một học sinh lên bảng viết hộ 3 từ khóa của phần vừa học.
\begin{itemize}
\item Không yêu cầu trình bày đẹp.
\item Cả lớp có thêm một điểm nhìn mới.
\item Thầy cô chỉ cần chốt lại bằng một câu ngắn.
\end{itemize}
\end{goiydungngay}
\begin{cauchot}
Đổi người cầm phấn đôi khi đủ để giữ cả tiết học tiếp tục.
\end{cauchot}
\ngancach

\section{Đổi kiểu trả lời từ nói sang chọn}
\begin{tinhhuong}{Tình huống}
Lớp mệt nên ngại phát biểu dài, nhưng vẫn cần tham gia để không trôi.
\end{tinhhuong}
\begin{goiydungngay}
Cho lớp chọn nhanh giữa hai hoặc ba phương án rồi mới mời giải thích 1 ý.
\begin{itemize}
\item Chọn trước giúp nhiều em dám nhập cuộc hơn.
\item Sau lựa chọn, chỉ cần 1 đến 2 bạn nói lý do.
\item Tiết học đổi nhịp mà không đổi mục tiêu nội dung.
\end{itemize}
\end{goiydungngay}
\begin{cauchot}
Khi lớp ngại nói dài, một lựa chọn nhanh là đường vòng rất hiệu quả vào sự tham gia.
\end{cauchot}
